\begin{definition}
\label{definition-pseudo-coherent}
Let $R$ be a ring. Denote $D(R)$ its derived category.
Let $m \in \mathbf{Z}$.
\begin{enumerate}
\item An object $K^\bullet$ of $D(R)$ is {\it $m$-pseudo-coherent}
if there exists a bounded complex $E^\bullet$ of finite free $R$-modules
and a morphism $\alpha : E^\bullet \to K^\bullet$ such that
$H^i(\alpha)$ is an isomorphism for $i > m$ and $H^m(\alpha)$
is surjective.
\item An object $K^\bullet$ of $D(R)$ is {\it pseudo-coherent}
if it is quasi-isomorphic to a bounded above complex of finite
free $R$-modules.
\item An $R$-module $M$ is called {\it $m$-pseudo-coherent}
if $M[0]$ is an $m$-pseudo-coherent object of $D(R)$.
\item An $R$-module $M$ is called
{\it pseudo-coherent}\footnote{This clashes with what is meant by
a pseudo-coherent module in \cite{Bourbaki-CA}.}
if $M[0]$ is a pseudo-coherent object of $D(R)$.
\end{enumerate}
\end{definition}